\chapter{Hardware Overview}
\label{chap:fpgalix-hw-overview}

\section{Overview}
\label{sec:fpgalix-overview}

This chapter illustrates the hardware realized for \prjname{FPGAlix} piece by piece, from the \gterm{ov7670}{OV7670} camera interface to the complete system.\footnote{Figure~\ref{fig:fpgalix-system-overview}, \repopath{soc_system.qsys} block diagram; also viewable directly through \gterm{platformdesigner}{Platform Designer} in \term{Quartus}, by opening \repopath{hw/quartus/soc_system.qsys}.}

As described in the \nameref{chap:preface}, the \gterm{fpga}{FPGA} fabric in \prjname{FPGAlix} is used to acquire the video stream, realizing a hardware interface that carries the data through the \gterm{msgdma}{mSGDMA} \term{IP} \cite{EP-IPUG}, an \term{Avalon Memory-Mapped} \gterm{dma}{DMA} engine capable of reading from an \gterm{avalonst}{Avalon-ST}\footnote{\glsentrydesc{avalonst}} stream and writing it to memory autonomously.

The remaining design question concerns how a byte-wide, framed data stream, subject to \gterm{backpressure}{backpressure} and originating in the camera's own clock domain, should reach the \gterm{msgdma}{mSGDMA}. \gterm{avalonst}{Avalon-ST}~\cite{AVALON-SPEC} constitutes the natural choice for this purpose: it is the standard \term{Intel}/\gterm{altera}{Altera} interconnect defined for unidirectional streaming transfers, providing a ready/valid handshake for backpressure and \gterm{sop}{SOP}/\gterm{sop}{EOP} (Start-of-packet/End-of-packet) markers for frame delimitation as part of the specification itself. Moreover, it is the interface natively supported by every other \term{IP} core involved in the pipeline (\hw{altera_avalon_dc_fifo} for the clock-domain crossing, \hw{altera_avalon_st_adapter} for width conversion, \hw{altera_avalon_sc_fifo} for buffering, and the \term{ST-to-MM} engine internal to \hw{altera_msgdma}). Implementing the camera capture core (\hw{ov7670_data_interface}) as an \gterm{avalonst}{Avalon-ST} source, rather than as a custom bus, allows it to be integrated directly into this chain without additional glue logic.

\section{The OV7670 camera and its adapter board}

\subsection{OV7670 overview}

The \gterm{ov7670}{OV7670} is a \term{CMOS} \gterm{vga}{VGA}\footnote{\glsentrydesc{vga}} image sensor that does not generate its own internal timing \cite{OV7670-DS}: correct operation requires an external master clock (\term{XCLK}), from which all subsequent timing is derived. Once clocked and configured, the sensor exposes the following signals:

\begin{itemize}
    \item \gterm{pclk}{PCLK}: pixel clock, generated by the sensor, synchronous with the data bus.
    \item \gterm{href}{HREF}: asserted during each valid line.
    \item \gterm{vsync}{VSYNC}: frame boundary marker.
    \item \term{D[7:0]}: 8-bit parallel data bus (format configurable as \gterm{bayermosaic}{RAW Bayer}, \gterm{yuv422}{YUV422}\footnote{\glsentrydesc{yuv422}}, or \gterm{rgb565}{RGB565}\footnote{\glsentrydesc{rgb565}}).
    \item \term{SIOC}/\term{SIOD}: \gterm{i2c}{SCCB}, a two-wire control interface derived from \gterm{i2c}{I\textsuperscript{2}C}, used to access the sensor's configuration registers.
    \item \term{RESET\#}, \term{PWDN}: reset and power-down control inputs.
\end{itemize}

\subsection{Rationale for a custom adapter board}

Direct connection of an \gterm{ov7670}{OV7670} breakout to the \term{DE1-SoC}'s \term{GPIO\_0} header by means of unshielded jumper wires does not provide adequate signal integrity at the sensor's clock frequency, up to 24 MHz for \gterm{pclk}{PCLK} at 640$\times$480 30 fps (12 MHz in \gterm{bayermosaic}{Bayer} mode)~\cite[Table 3-3]{OV7670-IG}: the resulting ringing and crosstalk are sufficient to corrupt sampling of \gterm{href}{HREF}, \gterm{vsync}{VSYNC}, and the data bus. To address this, a custom adapter was designed\footnote{See Appendix~\ref{app:camera-adapter-board} for the full schematics, bill of materials and photographs.}: two purpose-built \term{PCBs}, connected by a standard 40-pin \term{GPIO} ribbon cable (sources under \repopath{hw/camera_adapter/}):

\begin{itemize}
    \item \textbf{Transmitter (camera side):} mounted adjacent to the \gterm{ov7670}{OV7670} module. Each video signal (\gterm{pclk}{PCLK}, \gterm{href}{HREF}, \gterm{vsync}{VSYNC}, \term{D[7:0]}) is driven through an \term{SN74ALVC16244A} line buffer and onto the ribbon cable through a 33 $\Omega$ series resistor; \term{SDA} and \term{SCL} each get their own 4.7 k$\Omega$ pull-up resistor, one per line, two on this board.
    \item \textbf{Receiver (FPGA side):} mounted at the \term{DE1-SoC} \term{GPIO\_0} header. Each incoming video signal is terminated by a pair of 220 $\Omega$ resistors to \term{+3.3V} and to \term{DGND}; \term{SDA} and \term{SCL} are, likewise, each given their own 4.7 k$\Omega$ pull-up resistor, one per line, two on this board as well.
\end{itemize}

\subsection{Camera master clock generation}
\label{subsec:camera-master-clock}

The 24 MHz master clock (\term{Y1}, an \term{ECS-7050MV-240-BN-TR} \term{HCMOS} oscillator) is located on the transmitter board, adjacent to the sensor, and drives the sensor's \term{XCLK} directly through a 33 $\Omega$ series resistor; this clock is not routed across the ribbon cable. The alternative design, generating the clock on the \gterm{fpga}{FPGA}/receiver side and distributing it to the camera over the cable, would require an additional buffer stage on that side as well. Locating a dedicated cheap oscillator next to the sensor avoids the cost of this pricier second buffer and removes the signal-integrity risk associated with distributing a 24 MHz clock over the length of the ribbon cable.

\subsection{Line termination}

Two termination strategies are employed, one at each end of the cable:

\begin{itemize}
    \item \textbf{Transmitter side:} 33 $\Omega$ series (source) resistors on every buffered output. Source termination matches the buffer's low output impedance to the line, damping reflections launched at the driving end.
    \item \textbf{Receiver side:} each video signal is terminated by a Thevenin (split) termination: a 220 $\Omega$ resistor to \term{+3.3V} in parallel with a 220 $\Omega$ resistor to \term{DGND}. Their parallel combination presents 220 $\Omega$ $\parallel$ 220 $\Omega$ = 110 $\Omega$ to the line, matching the ribbon cable's estimated $\sim$110 $\Omega$ characteristic impedance. This absorbs reflections arriving at the receiving end and, as a side effect, biases each line to mid-rail whenever it is not actively driven.
\end{itemize}

\subsection{Video signal timing}

The two figures from the data-sheet \cite{OV7670-DS} below describe how the video pins behave and are the basis for the acquisition core described in Section~\ref{sec:fpgalix-acquisition-core}.

\begin{figure}[htbp]
    \centering
    \includegraphics[width=\linewidth]{OV7670_HTiming}
    \caption{OV7670 horizontal timing (source: \cite{OV7670-DS}).}
    \label{fig:ov7670-htiming}
\end{figure}

Within a frame, \gterm{href}{HREF}'s rising edge marks the start of a line and its falling edge marks the end. \term{D[7:0]} changes around \gterm{pclk}{PCLK}'s falling edge and is valid to be sampled on the following rising edge, one byte per \gterm{pclk}{PCLK} cycle for as long as \gterm{href}{HREF} stays high.


The number of bytes per pixel depends on the output format configured over \gterm{i2c}{SCCB}. \gterm{rgb565}{RGB565}, \gterm{rgb555}{RGB555}\footnote{\glsentrydesc{rgb555}}, \gterm{rgb444}{RGB444}\footnote{\glsentrydesc{rgb444}}, and \gterm{yuv422}{YUV422} transmit each pixel as 2 bytes over 2 \gterm{pclk}{PCLK} cycles. \gterm{bayermosaic}{RAW Bayer} and processed \gterm{bayermosaic}{Bayer} transmit each pixel as a single byte, and, to compensate, the camera halves \gterm{pclk}{PCLK}'s frequency in that mode, so the pixel period (and with it, line and frame timing, Figure~\ref{fig:ov7670-vtiming} below) stays the same regardless of format.

Byte ordering and bit layout within each format are given in the data-sheet \cite{OV7670-DS}, not reproduced here.

\begin{figure}[htbp]
    \centering
    \includegraphics[width=\linewidth]{OV7670_VGAFrameTiming}
    \caption{OV7670 VGA frame timing (source: \cite{OV7670-DS}).}
    \label{fig:ov7670-vtiming}
\end{figure}

Figure~\ref{fig:ov7670-vtiming} gives the same picture at the scale of a whole frame. A frame begins on \gterm{vsync}{VSYNC}'s falling edge; after a blanking gap, the first line begins at the first \gterm{href}{HREF} rising edge (row 0). The last active row (row 479) ends on its \gterm{href}{HREF} falling edge, and after a further blanking gap, \gterm{vsync}{VSYNC}'s next rising edge signals that the frame is complete. Over a full frame this amounts to 480 active rows of 640 pixels each.

\term{HSYNC} appears in Figure~\ref{fig:ov7670-vtiming} for reference only: it is not a physical output pin on the \gterm{ov7670}{OV7670}. Only \gterm{pclk}{PCLK}, \gterm{href}{HREF}, and \gterm{vsync}{VSYNC} are, and they are what the custom acquisition \term{IP} core actually uses to delimit frames and lines.

\longsection{Generating and resetting \term{PCLK}}{\texttt{pclk\_reset\_controller}}
\label{sec:fpgalix-pclk-reset}

Unlike \hw{sys_clk} (\hw{clk_0}, 50 MHz, free-running from the board oscillator), \gterm{pclk}{PCLK} does not exist until the camera is powered up and out of reset: it is generated by the \gterm{ov7670}{OV7670} itself, derived internally from the clock generator on the transmitter board\footnote{See Section~\ref{subsec:camera-master-clock}.}, and only starts toggling once the sensor is running.

\hw{pclk_reset_controller} is the custom \gterm{platformdesigner}{Platform Designer} component that turns this externally generated signal into a proper clock domain for the rest of the design and manages the two resets that go with it.

\begin{figure}[htbp]
    \centering
    \includesvg[width=\linewidth]{PCLKDistribution}
    \caption{Distribution of \term{PCLK} and its reset, from the camera to every \term{pclk}-domain block.}
    \label{fig:pclk-distribution}
\end{figure}

\subsection{\term{PCLK} as a \term{Platform Designer} clock}

\hw{pclk_reset_controller} takes the physical \gterm{pclk}{PCLK} pin as \hw{pclk_in} and re-exports it as \hw{pclk_out}, a genuine \gterm{platformdesigner}{Platform Designer} clock source: \ccode{pclk_out <= pclk_in} in the \gterm{vhdl}{VHDL}, a pattern \term{Quartus} recognizes and implements with a global clock buffer rather than ordinary fabric routing.

\subsection{The reset problem}

The domain's reset is controlled by a single \hw{sw_release} register bit, reachable from the main \gterm{avalonmm}{Avalon-MM}\footnote{\glsentrydesc{avalonmm}} bus running in the \hw{sys_clk} domain. A synchronizer for it cannot be built the ordinary way: logic sitting in the \gterm{pclk}{PCLK} domain normally needs \gterm{pclk}{PCLK} edges to deassert its reset cleanly, but \gterm{pclk}{PCLK} itself does not exist while the camera is held in reset, so waiting on it to release its own gating reset would deadlock.

\hw{pclk_reset_controller} resolves this by deriving two separate reset outputs from the same condition, \hw{sys_reset_n} and \hw{sw_release}:

\begin{itemize}
    \item \hw{camera_reset_n}: purely combinatorial, entirely in the \hw{sys_clk} domain, no dependency on \gterm{pclk}{PCLK} at all. This drives the camera's physical reset pin directly, so it is what gets the camera running and \gterm{pclk}{PCLK} toggling again.
    \item \hw{pclk_reset_n}: a two-flip-flop synchronizer, clocked on \hw{pclk_in}, that brings the same condition into the \gterm{pclk}{PCLK} domain for the acquisition logic. Assert is asynchronous, both flip-flops clear immediately when the condition drops, so the domain stays in reset even if \gterm{pclk}{PCLK} stops ticking mid-frame, while deassert is synchronous, two clean \gterm{pclk}{PCLK} edges after the camera comes back up, giving proper metastability protection for the moment \gterm{pclk}{PCLK} starts again.
\end{itemize}

Because \hw{camera_reset_n} has no dependency on \gterm{pclk}{PCLK}, releasing it is always safe regardless of whether \gterm{pclk}{PCLK} is running.

\subsection{Usage sequence}

The dependency above dictates the order software has to follow:

\begin{enumerate}
    \item Write 0 to the \hw{sw_release} register: both \hw{camera_reset_n} and \hw{pclk_reset_n} are held low.
    \item Write 1 to the \hw{sw_release} register: \hw{camera_reset_n} rises immediately, releasing the camera; \hw{pclk_reset_n} follows two \gterm{pclk}{PCLK} edges later, once the camera starts driving \gterm{pclk}{PCLK} again.
    \item Only once the camera is out of reset can it answer over \gterm{i2c}{SCCB}, so its configuration has to come after this step, not before it.
\end{enumerate}

\longsection{The camera acquisition core}{\texttt{ov7670\_data\_interface}}
\label{sec:fpgalix-acquisition-core}

\hw{ov7670_data_interface} is a custom \gterm{platformdesigner}{Platform Designer} component that turns the camera's \gterm{pclk}{PCLK}/\gterm{href}{HREF}/\gterm{vsync}{VSYNC}/\term{D[7:0]} signals into the \gterm{avalonst}{Avalon-ST} stream the rest of the pipeline expects. Both its \gterm{avalonst}{Avalon-ST} source and its \gterm{avalonmm}{Avalon-MM} \gterm{registerinterface}{register interface} run in the \gterm{pclk}{PCLK} domain, the camera pixel clock. A single control register (\hw{ctrl}) lets software enable acquisition (\hw{ctrl_enabled} bit), drive the camera's power-down pin (\hw{ctrl_cam_pwdn} bit), and clear the dropped-frame counter (\hw{ctrl_clear_counters} bit); a second, read-only register (\hw{stat_dropped_cnt}) counts frames the core had to discard.

Because this register interface resides in the \gterm{pclk}{PCLK} domain rather than \hw{sys_clk}, software cannot reach it directly from the \gterm{lwhpsbridge}{LW HPS} bridge\footnote{\glsentrydesc{lwhpsbridge}}: \hw{mm_clock_crossing_bridge_0}, an \gterm{avalonmm}{Avalon-MM} clock-crossing bridge, is placed between the two domains for exactly this purpose, translating an \gterm{avalonmm}{Avalon-MM} access from \hw{sys_clk} into \gterm{pclk}{PCLK} and back. It exists solely to give \hw{ov7670_data_interface_0}'s register file a path from software; no other component in the design needs it.

Functionally, the core forwards each active pixel byte onto the \gterm{avalonst}{Avalon-ST} bus as it arrives from the camera, treating an entire camera frame as a single \gterm{avalonst}{Avalon-ST} packet: the \gterm{sop}{SOP} signal marks its first byte, the \gterm{sop}{EOP} one its last. As long as the downstream side keeps accepting data, this packet runs uninterrupted for the whole frame. If a backpressure event arrives instead, that is, the downstream side stops accepting data before the frame is complete, the core does not wait for it to recover: it abandons the frame in progress (drop), and inserts the \gterm{sop}{EOP} at the first opportunity, once the downstream side accepts data again, in place of the remaining pixel data. \gterm{msgdma}{mSGDMA} then terminates that transfer early, and software can tell the frame is incomplete from the number of bytes actually transferred, rather than being handed truncated pixel data as if it were a valid frame.

This behavior is implemented as a three-state \gterm{fsm}{machine}, built directly on the \gterm{href}{HREF}/\gterm{vsync}{VSYNC} signaling described above and shown in Figure~\ref{fig:ov7670-vtiming}:

\begin{figure}[htbp]
    \centering
    \includesvg[width=\linewidth]{CameraDataInterfaceFSM}
    \caption{State machine of \term{ov7670\_data\_interface}.}
    \label{fig:ov7670-data-interface-fsm}
\end{figure}

\begin{itemize}
    \item \textbf{WAIT\_SYNC}: the reset and idle state. It waits for \gterm{vsync}{VSYNC}'s falling edge, the frame-start marker from Figure~\ref{fig:ov7670-vtiming}, with acquisition enabled (\hw{ctrl_enabled} = 1), and moves to ACQUIRE.
    \item \textbf{ACQUIRE}: for every \gterm{pclk}{PCLK} cycle with \gterm{href}{HREF} asserted, one byte from the \term{D[7:0]} data line is transferred onto the \gterm{avalonst}{Avalon-ST} bus. \gterm{sop}{SOP} is raised on the very first byte and \gterm{sop}{EOP} on the \ccode{FRAME_SIZE}-th byte (\ccode{FRAME_SIZE}, a \gterm{vhdl}{VHDL} generic constant, being a count of bytes, not pixels, see Figure~\ref{fig:ov7670-vtiming} and Figure~\ref{fig:ov7670-htiming}), closing the frame and returning to the WAIT\_SYNC status: the byte count, not the next \gterm{vsync}{VSYNC} edge, is what ends the packet. If the downstream side deasserts the \term{READY} signal at any point (backpressure), the core does not stall; it counts the frame as dropped and moves to WAIT\_READY instead.
    \item \textbf{WAIT\_READY}: reached after a drop. The core waits for \term{READY} to return, then emits a bare \gterm{sop}{EOP}, so the downstream FIFO can discard the incomplete packet cleanly, and goes back to WAIT\_SYNC. Each further \gterm{vsync}{VSYNC} falling edge seen while still waiting counts as another dropped frame, so a downstream stall spanning several frames is reported accurately rather than silently
swallowed.
\end{itemize}

\longsection{The Avalon-ST pipeline}{from \texttt{ov7670\_data\_interface} to \term{mSGDMA}}

Section~\ref{sec:fpgalix-overview} established the rationale for streaming pixel data from the camera to \gterm{msgdma}{mSGDMA} over \gterm{avalonst}{Avalon-ST}, such that the \gterm{hps}{HPS} is limited to configuring a transfer and being notified upon its completion. This chapter examines that stream in detail, from the camera acquisition \term{IP}'s \gterm{avalonst}{Avalon-ST} source port (\hw{ov7670_data_interface_0.st_source}) to the write into \gterm{hps}{HPS} \term{DDR3}; the corresponding block diagram is Figure~\ref{fig:fpgalix-system-overview}.

\subsection{The pipeline}

The following \term{IP} instances, shown in Figure~\ref{fig:fpgalix-system-overview}, are traversed in order: \hw{ov7670_data_interface_0}, \hw{dc_fifo_0}, \hw{st_adapter_0}, \hw{sc_fifo_0}, and \hw{msgdma_0}. Of these components, only \hw{dc_fifo_0} spans a clock-domain boundary: its \hw{in_clk} input port is driven by \hw{pclk_out} from \hw{pclk_reset_controller_0}, while its \hw{out_clk} output port is driven by \hw{sys_clk}. Every component downstream of \hw{dc_fifo_0}, namely \hw{st_adapter_0}, \hw{sc_fifo_0}, and \hw{msgdma_0}, operates entirely within the \hw{sys_clk} domain.

\subsection{Clock-domain crossing: \texttt{dc\_fifo\_0}}

The clock domain crossing is performed at this point, immediately following the camera acquisition core and prior to any width conversion, for two reasons. First, crossing an 8-bit stream requires substantially less synchronization logic than crossing a wider stream, since a dual-clock FIFO requires synchronizers proportional to its data width. Second, this arrangement confines the dependency on \hw{pclk}, whose frequency varies with the camera's output format and which ceases entirely while the camera is held in reset, to the smallest possible portion of the design; every subsequent stage operates against a single, stable clock.

\hw{dc_fifo_0} is configured with a depth of 65536 entries at 8 bits, providing an elastic buffer sufficient to absorb the rate difference between \hw{pclk} and \hw{sys_clk}, as well as transient stalls on the \hw{sys_clk} side, without requiring backpressure to propagate to the camera interface. Backpressure that does reach \hw{ov7670_data_interface} is instead handled by it through discarding the affected frame rather than stalling.

\subsection{Width conversion: \texttt{st\_adapter\_0}}
\label{subsec:fpgalix-st-adapter}

\hw{st_adapter_0} converts the stream from 8-bit, single-\gterm{symbol}{symbol} \gterm{beat}{beats} to 128-bit, 16-symbol beats. Performing this conversion after the clock-domain crossing, rather than before it, keeps the synchronization logic in \hw{dc_fifo_0} narrow while providing \hw{msgdma_0} and its destination, \gterm{hps}{HPS} \term{DDR3}, with the wide beats their burst mechanisms are designed around: a 128-bit \gterm{avalonmm}{Avalon-MM} write makes substantially more efficient use of \term{SDRAM} burst bandwidth than the equivalent data transferred one byte at a time.

\subsection{Buffering: \texttt{sc\_fifo\_0}}
\label{subsec:fpgalix-sc-fifo}

A single-clock, 128-bit buffer \term{FIFO}, 8192 entries deep, is placed entirely within the \hw{sys_clk} domain between the width converter and \gterm{msgdma}{mSGDMA}. It accommodates timing variation between the arrival of data from the pipeline and the availability of \gterm{msgdma}{mSGDMA} to issue a burst, for example while \gterm{msgdma}{mSGDMA} is between bursts or the \gterm{hps}{HPS} memory bus is occupied by other traffic.

\subsection{\term{mSGDMA}: from \term{Avalon-ST} to \term{DDR3}}
\label{subsec:fpgalix-msgdma}

The \gterm{msgdma}{mSGDMA} is Intel's configurable \gterm{dma}{DMA} \term{IP} core, moving data between two memory spaces (i.e.\ Avalon-MM to Avalon-MM), or between a streaming interface and memory (i.e.\ Avalon-MM to Avalon-ST or vice versa), under the control of preloaded instructions called \emph{descriptors}.

Internally, the core is built around one central module, the \emph{dispatcher}, to which up to two further modules can be attached: a \emph{Read Host}, used when the source of the transfer is memory-mapped, and a \emph{Write Host}, used when its destination is memory-mapped; when both ends of the transfer are memory-mapped, both are attached. The dispatcher receives the descriptors from software, decodes them one at a time, and issues the corresponding read and/or write commands to whichever of the two is present, collecting their response to report the transfer's outcome back to software.

Each descriptor specifies, at minimum, the source and destination addresses and the number of bytes to transfer. A sequence of descriptors queued in advance is what allows the dispatcher to execute successive transfers without further software intervention between one and the next.

In \prjname{FPGAlix}, the \gterm{msgdma}{mSGDMA}, instantiated as \hw{msgdma_0}, is used in its Streaming-to-Memory-Mapped variant, since its source, \hw{sc_fifo_0}, is an \gterm{avalonst}{Avalon-ST} interface, while its final destination is the \term{DDR3} memory: only a Write Host is therefore attached to the dispatcher, as shown in Figure~\ref{fig:fpgalix-msgdma-st-to-mm}.

\begin{figure}[htbp]
    \centering
    \includegraphics[width=\linewidth]{MsgdmaStToMmIntelManual}
    \caption{mSGDMA module configuration with a Write Host attached, supporting \term{Avalon-ST} to \term{Avalon-MM} transfers (source: \cite[Figure 105, p.~360]{EP-IPUG}).}
    \label{fig:fpgalix-msgdma-st-to-mm}
\end{figure}

The Write Host receives \hw{sc_fifo_0}'s output on its \gterm{avalonst}{Avalon-ST} sink port and issues the corresponding writes out of its own \gterm{avalonmm}{Avalon-MM} master port with the width set to 128 bit. Bursting is enabled on this instance, with a maximum burst count of 512 beats, for bursts of up to 8~\term{KiB}.

Software drives \gterm{msgdma}{mSGDMA} through three interfaces reached over the \gterm{lwhpsbridge}{LW HPS} bridge: the \gterm{csr}{CSR} (\hw{csr}), \hw{descriptor_slave}, and \ccode{response}. A descriptor's fields are written directly into \hw{descriptor_slave}, staged until its \term{Go} bit commits it into the core's internal descriptor \term{FIFO}; queued descriptors are then processed autonomously, with completion reported both through an interrupt and a status write to \ccode{response}.

Packet support is enabled on this instance specifically, so that the transfer ends as soon as \gterm{msgdma}{mSGDMA} observes an incoming \gterm{sop}{EOP} on its sink, rather than only once the descriptor's full length has been transferred; the length field thus becomes an upper bound rather than an exact size. This is precisely why \hw{ov7670_data_interface} frames the stream with \gterm{sop}{SOP} and \gterm{sop}{EOP} in the first place (Section~\ref{sec:fpgalix-acquisition-core}): a complete frame reaches \ccode{FRAME_SIZE} bytes before its own \gterm{sop}{EOP}, while a dropped frame's bare \gterm{sop}{EOP} arrives early. The \ccode{response} port reports the actual byte count transferred for each descriptor, which is how software tells a short, dropped frame apart from a complete one.

\subsection{Why a dedicated \term{SDRAM} port: bypassing the \term{L3} interconnect}
\label{sec:fpgalix-sdram-port}

\begin{figure}[htbp]
    \centering
    \includegraphics[width=\linewidth]{HPSBlockDiagramIntelManual}
    \caption{\term{HPS} block diagram, annotated to highlight the two candidate paths from the \term{FPGA} fabric to \term{DDR3}: via the general-purpose bridge and the \term{L3} Main Switch (yellow), or via a dedicated \term{SDRAM} port (green) (source: \cite[Figure 2-2, p.~2-4]{CV-TRM}, annotated).}
    \label{fig:fpgalix-hps-block-diagram}
\end{figure}

Reaching \term{DDR3} from the \gterm{fpga}{FPGA} fabric is not a single obvious path: Figure~\ref{fig:fpgalix-hps-block-diagram} highlights the two options available. One attaches to the general-purpose \gterm{fpga}{FPGA}-to-\gterm{hps}{HPS} bridge already exposed by the \gterm{hps}{HPS} (yellow path): a master issues transactions into the bridge, which lands on the \term{L3} Main Switch, the same switch that also carries traffic from all built-in \gterm{hps}{HPS} peripherals (like the Ethernet controller and the \term{SD} card), before finally reaching the \term{SDRAM} Controller Subsystem. The other exposes one of the \term{SDRAM} Controller Subsystem's own dedicated ports directly to the \gterm{fpga}{FPGA} fabric (green path, the ``1-6 Masters'' in the diagram), bypassing the \term{L3} Main Switch entirely.

Given the camera's throughput (640$\times$480 at 30 fps, 2 bytes per pixel, i.e.\ roughly 17.6 \term{MiB}/s) and the fabric resources available, \prjname{FPGAlix} takes the second path: a dedicated port (\hw{f2h_sdram1}), configured in \gterm{platformdesigner}{Platform Designer}, exposes the \term{SDRAM} Controller Subsystem directly to the \gterm{fpga}{FPGA} fabric, 128 bit wide.

\begin{figure}[htbp]
    \centering
    \includegraphics[width=\linewidth]{HPSRamCtrlSubsycFromIntelManual}
    \caption{\term{SDRAM} Controller Subsystem high-level block diagram: the Multi-Port Front End arbitrates the \term{MPU}, \term{L3}, and \term{FPGA} fabric interfaces down to the single-port controller. Traffic entering through the \term{FPGA} fabric interface bypasses the \term{MPU} subsystem entirely, including its \term{L2} cache, so it never becomes visible to the \term{CPU}'s cache system (source: \cite[Figure 12-1, p.~12-2]{CV-TRM}).}
    \label{fig:fpgalix-hps-sdram}
\end{figure}

This works because the \term{SDRAM} controller (Figure~\ref{fig:fpgalix-hps-sdram}) exposes a Multi-Port Front End (\gterm{mpfe}{MPFE}) that arbitrates a handful of incoming ports down to the single port that actually talks to the memory chips. Two of them are already taken: one feeds the \gterm{mpu}{MPU}, at 64 bit bus width, the other the \term{L3} interconnect, at 32 bit bus width. The rest are left available to the \gterm{fpga}{FPGA} fabric, and this project uses two of them, as already seen: \hw{f2h_sdram1} for the video stream, through \hw{msgdma_0}, and \hw{f2h_sdram0} for \gterm{jtag}{JTAG} debugging access, through \hw{f2sdram_only_master}.

\longsection{The 7-segment displays}{\texttt{hex\_display\_controller}}
\label{sec:fpgalix-hex-display}

The \term{DE1-SoC}'s six 7-segment displays are driven by \hw{hex_display_controller}, a custom \gterm{platformdesigner}{Platform Designer} component. A single register controls all six displays together, and any six-digit decimal number written to it is displayed accordingly.

\subsection{Register interface}
\label{sec:fpgalix-hex-register}

The component exposes a single 32-bit \gterm{avalonmm}{Avalon-MM} register on the \gterm{lwhpsbridge}{LW HPS}-to-\gterm{fpga}{FPGA} bridge, with the following layout:

\begin{itemize}
    \item bit 31: \ccode{enabled} (1 displays the value below, 0 blanks every digit)
    \item bits [19:0]: the value to be displayed, in the range 0...999999
    \item any value above 999999 in bits [19:0] forces every digit to a dash, provided bit 31 is set; if bit 31 is clear, every digit is blanked regardless of this value
\end{itemize}

\subsection{From binary to seven segments}

The 20-bit binary value is converted to six \term{BCD} digits using the \term{Double Dabble} algorithm~\cite{PEDRONI-VHDL}; each digit then is driven onto one of the six displays. Leading zeroes are suppressed: only the digits needed to represent the value are shown, with the units digit always shown even when the value is zero.

\section{System overview}
\label{sec:fpgalix-system-overview}

The diagram in Figure~\ref{fig:fpgalix-system-overview} is the complete \gterm{platformdesigner}{Platform Designer} system (\repopath{soc_system.qsys}), placing the blocks discussed in the preceding chapters within the wider system: the \gterm{hps}{HPS}, its lightweight and full bridges into the \gterm{fpga}{FPGA} fabric, and the \term{DE1-SoC} peripherals wired directly to dedicated \gterm{hps}{HPS} pins.

Several blocks shown here are not covered by a dedicated chapter: the \term{LED}, \term{switch}, and \term{key} \term{PIOs}, the \term{System ID} core, the \term{JTAG UART}, the \term{Interrupt Latency Counter}, and the \gterm{hps}{HPS} peripherals wired to their own dedicated pins (\term{Ethernet}, \term{USB}, \term{SD/MMC}, \term{QSPI} flash, and the various serial buses). None of these originate from this project's own design work; they were inherited and combined from the \term{Cyclone V} \term{Golden Hardware Reference Design} \cite{GHRD} and from the \term{DE1-SoC} \term{Embedded Linux} Systems guide \cite{UPM-DE1}, the two starting points this system was built from, and are documented no further here. The chapters above account for everything specific to this project: the camera, its clocking, the path from acquisition to \term{DDR3}, and the 7-segment displays.

\begin{landscape}
\begin{figure}[p]
    \centering
    \includesvg[width=\linewidth,height=\textwidth,keepaspectratio]{FPGAlixQsysDiagram}
    \caption{\texttt{soc\_system.qsys} block diagram.}
    \label{fig:fpgalix-system-overview}
\end{figure}
\end{landscape}

\section{Resource utilization}
\label{sec:fpgalix-resource-utilization}

Table~\ref{tab:fpgalix-resource-utilization} reports the \gterm{fpga}{FPGA} fabric resources consumed by the complete system of Figure~\ref{fig:fpgalix-system-overview}, as reported by \gterm{quartusprime}{Quartus Prime}'s fitter after place-and-route, on the \term{DE1-SoC}'s \term{Cyclone V} \term{5CSEMA5F31C6} device.

\begin{table}[htbp]
\centering
\begin{tabularx}{\textwidth}{@{} l X X X @{}}
\toprule
\textbf{Resource} & \textbf{Used} & \textbf{Available} & \textbf{Utilization} \\
\midrule
\term{ALM}s & 4{,}725 & 32{,}070 & 15\% \\
\lightmidrule
Total registers & 7{,}434 & --- & --- \\
\lightmidrule
Total block memory bits & 2{,}864{,}980 & 4{,}065{,}280 & 70\% \\
\lightmidrule
\term{DSP} blocks & 0 & 87 & 0\% \\
\lightmidrule
\term{I/O} pins & 233 & 457 & 51\% \\
\bottomrule
\end{tabularx}
\caption{\term{FPGA} resource utilization of the complete \prjname{FPGAlix} system (source: \term{Quartus Prime} fitter report).}
\label{tab:fpgalix-resource-utilization}
\end{table}

Logic and \term{DSP} usage are unremarkable, at 15\% and 0\% respectively. Block memory, at 70\%, is the one figure worth commenting on, and follows directly from the two capture-path \term{FIFO}s sized generously to mitigate the backpressure problems discussed in \nameref{chap:fpgalix-vision-core}: \hw{dc_fifo_0} (65536 entries at 8 bits) and \hw{sc_fifo_0} (8192 entries at 128 bits) together account for most of it.

\section{Platform Designer project screenshots}
\label{sec:fpgalix-platformdesigner-screenshots}

Figures~\ref{fig:fpgalix-qsys-connections-1}--\ref{fig:fpgalix-qsys-connections-3} show the component and connection listing of \repopath{hw/quartus/soc_system.qsys} directly from \gterm{platformdesigner}{Platform Designer}, for reference alongside the block diagram of Figure~\ref{fig:fpgalix-system-overview}.

\begin{landscape}
\begin{figure}[p]
    \centering
    \includegraphics[width=\linewidth,height=\textwidth,keepaspectratio]{FPGAlixQsys_1}
    \caption{\texttt{soc\_system.qsys} component connections in \term{Platform Designer}, part 1 of 3.}
    \label{fig:fpgalix-qsys-connections-1}
\end{figure}
\end{landscape}

\begin{landscape}
\begin{figure}[p]
    \centering
    \includegraphics[width=\linewidth,height=\textwidth,keepaspectratio]{FPGAlixQsys_2}
    \caption{\texttt{soc\_system.qsys} component connections in \term{Platform Designer}, part 2 of 3.}
    \label{fig:fpgalix-qsys-connections-2}
\end{figure}
\end{landscape}

\begin{landscape}
\begin{figure}[p]
    \centering
    \includegraphics[width=\linewidth,height=\textwidth,keepaspectratio]{FPGAlixQsys_3}
    \caption{\texttt{soc\_system.qsys} component connections in \term{Platform Designer}, part 3 of 3.}
    \label{fig:fpgalix-qsys-connections-3}
\end{figure}
\end{landscape}